\documentclass[11pt]{article}
\usepackage[margin=1in]{geometry}
\usepackage{amsmath,amssymb,amsthm,mathtools}
\usepackage{graphicx}
\usepackage{microtype}
\usepackage[colorlinks=true,linkcolor=blue,citecolor=blue,urlcolor=blue]{hyperref}

\newtheorem{theorem}{Theorem}
\newtheorem{lemma}{Lemma}
\newcommand{\F}{\mathbb F_2^+}
\newcommand{\cH}{\mathcal H}
\newcommand{\D}{\mathbb D}
\renewcommand{\thesection}{\arabic{section}.}

\title{An exact fixed-chaos slice of the\\
$L^2$ Haar-multishift commutant}
\author{Candidate partial result for arXiv:2005.04648}
\date{13 August 2026}

\begin{document}
\maketitle

\noindent\fbox{\parbox{0.94\linewidth}{\textbf{Status.}
Candidate partial result, likely valid.  We obtain an exact norm and
boundedness criterion for a generator in any one fixed Haar-chaos order.
The unrestricted commutant problem is \emph{not} claimed.}}

\section{Source question and result}

Let $h$ be the first Haar function and $V_0,V_1$ the two dyadic dilation
operators.  For a mean-zero function $f$, Astashkin and Terekhin define
\[
 T_fV^\alpha h=V^\alpha f.
\]
The source notes that $T_f$ commutes with the Haar multishift and states that
a comprehensive description of its commutant is open even in $L^2$
\textup{[1]}.  It completely treats first-order Haar chaoses.  The excerpt
below records the open problem.

\begin{figure}[h]
  \centering
  \includegraphics[width=0.86\linewidth]{figures/open_problem_crop.png}
  \caption{Definition of $T_f$ and the open $L^2$ commutant problem on
  page 7 of arXiv:2005.04648.}
\end{figure}

We solve the boundedness problem when $f$ belongs to one prescribed Haar
chaos.  Use the source parametrization
\begin{equation}\label{eq:f}
 f=\sum_{k_1,\ldots,k_d\ge0}\xi_{k_1,\ldots,k_d}
 V_0^{k_1}V_1V_0^{k_2}\cdots V_1V_0^{k_d}h
\end{equation}
and its symbol
\[
 \widehat f(z_1,\ldots,z_d)
 =\sum_{k_1,\ldots,k_d\ge0}\xi_{k_1,\ldots,k_d}
 z_1^{k_1}\cdots z_d^{k_d}.
\]
Put $R=2^{-1/2}$, and use the convention $H^2(\D_R^0)=\mathbb C$.

\begin{theorem}[fixed-chaos commutant formula]\label{thm:main}
Let $d\ge1$ and let $f\in\cH_{ch}^d\cap L^2$ have the expansion
\eqref{eq:f}.  The densely defined operator $T_f$ extends boundedly to the
mean-zero subspace $L_0^2[0,1]$ if and only if
\begin{equation}\label{eq:criterion}
 \sup_{|w|<R}
 \bigl\|\widehat f(w,\,\cdot)\bigr\|_{H^2(\D_R^{d-1})}<\infty.
\end{equation}
In that case its exact norm is
\begin{equation}\label{eq:norm}
 \boxed{\quad
 \|T_f\|_{L_0^2\to L_0^2}
 =2^{-(d-1)/2}\sup_{|w|<R}
 \bigl\|\widehat f(w,\,\cdot)\bigr\|_{H^2(\D_R^{d-1})}.
 \quad}
\end{equation}
For $d=1$, this is exactly the scalar $H^\infty(\D_R)$ formula proved in
the source.
\end{theorem}

\section{Proof intuition}

In the normalized Haar basis, the two multishift generators are the left
creation operators on binary words, whereas every commuting $T_f$ appends
the coefficient word on the right.  A word in one fixed chaos order has a
distinguished initial run of zeros.  Removing that run leaves a terminal
word from a suffix-free family, so the corresponding right-creation ranges
are orthogonal.

The only remaining interaction is convolution along the initial zero run.
That interaction is an ordinary unilateral shift.  Its Wold model converts
$T_f$ into multiplication by an $\ell^2$-valued analytic function: the
$\ell^2$ coordinates are precisely the remaining $d-1$ chaos variables.
The multiplier norm is therefore an $H^\infty$ norm in the first variable
and an $H^2$ norm in all the others.  This orthogonality is lost as soon as
different chaos orders are mixed.

\section{Free-word reduction}

Let $\F$ be the set of finite binary words, including the empty word, with
concatenation denoted by juxtaposition.  The normalized functions
\[
 e_\alpha=2^{|\alpha|/2}V^\alpha h,\qquad \alpha\in\F,
\]
form an orthonormal basis of $L_0^2$.  If
$f=\sum_{\beta\in\F}a_\beta e_\beta$, then a direct calculation gives
\begin{equation}\label{eq:right}
 T_fe_\alpha=\sum_{\beta\in\F}a_\beta e_{\alpha\beta}.
\end{equation}
Thus the coefficient word is appended on the right.  Define the right
creation isometry $R_\beta e_\alpha=e_{\alpha\beta}$.

For $d>1$, every word occurring in \eqref{eq:f} has a unique factorization
\[
 \beta=0^{k_1}\gamma,
 \qquad
 \gamma=1,0^{k_2}\cdots1,0^{k_d}.
\]
Let $\Gamma_d$ be the family of these $\gamma$; put
$\Gamma_1=\{\varnothing\}$.  The family $\Gamma_d$ is suffix-free.  Indeed,
if $\gamma$ were a proper suffix of $\delta$, the deleted nonempty prefix
of $\delta$ would begin with $1$.  It would contribute an additional one,
although both words have exactly $d-1$ ones.  This is impossible.

Consequently the ranges of $R_\gamma$, $\gamma\in\Gamma_d$, are pairwise
orthogonal: a word ending in two distinct terminal words would force one
of them to be a suffix of the other.  Write
\[
 b_{k,\gamma}=a_{0^k\gamma},\qquad
 \phi_\gamma(z)=\sum_{k\ge0}b_{k,\gamma}z^k.
\]
Equation \eqref{eq:right} now becomes
\begin{equation}\label{eq:factor}
 T_f=\sum_{\gamma\in\Gamma_d}R_\gamma\phi_\gamma(R_0),
 \qquad
 \|T_fx\|_2^2
 =\sum_{\gamma\in\Gamma_d}\|\phi_\gamma(R_0)x\|_2^2.
\end{equation}

\section{Vector-valued Hardy multiplier}

We record the functional-analytic step, including the infinite sum.

\begin{lemma}\label{lem:column}
Let $S$ be a pure isometry on a Hilbert space, and let
$\Phi=(\phi_j)_{j\in J}$ be an $\ell^2(J)$-valued analytic function whose
Taylor coefficients are square summable.  On the algebraic Wold span set
\[
 Cx=(\phi_j(S)x)_{j\in J}.
\]
Then $C$ is bounded if and only if
$\Phi\in H^\infty(\D;\ell^2(J))$, and
\[
 \|C\|=\sup_{|z|<1}\Bigl(\sum_{j\in J}|\phi_j(z)|^2\Bigr)^{1/2}.
\]
\end{lemma}

\begin{proof}
The Wold unitary identifies $S$ with multiplication by $z$ on
$H^2(\D;\mathcal W)$, where $\mathcal W=\ker S^*$.  Under this unitary,
$C$ is pointwise multiplication
\[
 g(z)\longmapsto \Phi(z)\otimes g(z)
\]
from $H^2(\D;\mathcal W)$ to
$H^2(\D;\ell^2(J)\otimes\mathcal W)$.  Boundary integration gives the
upper estimate by $\|\Phi\|_{H^\infty}$, and normalized reproducing kernels
give the reverse estimate at every point of the disk.  For infinite $J$,
apply this argument first to finite coordinate projections.  Their squared
pointwise norms increase to $\|\Phi(z)\|_{\ell^2}^2$, which proves both
necessity and the asserted norm identity.
\end{proof}

The isometry $R_0$ in \eqref{eq:factor} is pure: its wandering space is
spanned by the vacuum and by words ending in $1$.  Orthogonality of the
$R_\gamma$ makes $T_f$ isometrically equivalent to the column in
Lemma~\ref{lem:column}.  Hence
\begin{equation}\label{eq:vector}
 \|T_f\|=sup_{|z|<1}
 \left(\sum_{\gamma\in\Gamma_d}|\phi_\gamma(z)|^2\right)^{1/2}.
\end{equation}

It remains only to translate the normalization.  For
$\gamma=1,0^{k_2}\cdots1,0^{k_d}$, the word in \eqref{eq:f} has length
$k_1+\cdots+k_d+d-1$, so
\[
 \phi_\gamma(z)
 =2^{-(d-1+k_2+\cdots+k_d)/2}
 \sum_{k_1\ge0}\xi_{k_1,\ldots,k_d}(z/\sqrt2)^{k_1}.
\]
Putting $w=z/\sqrt2$ in \eqref{eq:vector}, and using the coefficient norm
of $H^2(\D_R^{d-1})$, gives
\[
 \sum_{\gamma\in\Gamma_d}|\phi_\gamma(z)|^2
 =2^{-(d-1)}
 \|\widehat f(w,\,\cdot)\|_{H^2(\D_R^{d-1})}^2.
\]
This proves Theorem~\ref{thm:main}.

\small
\section{Strictness, computation, and limitation}

Condition \eqref{eq:criterion} is stronger than $f\in L^2$.  For example,
fix all tail indices $k_2=\cdots=k_d=0$ and choose the first-variable symbol
\[
 g(w)=(1-\sqrt2w)^{-1/4}.
\]
Then $g\in H^2(\D_R)$ but $g\notin H^\infty(\D_R)$.  The corresponding
order-$d$ chaos belongs to $L^2$, while Theorem~\ref{thm:main} says that
$T_f$ is unbounded.  At the opposite normalization extreme,
$f=V_1h$ gives $\|T_f\|=2^{-1/2}$, as it must.

The accompanying finite-word audit uses an order-three polynomial with four
terminal suffixes.  The normalized vector norm and independently scaled
mixed norm agree to twelve decimal places, while exact compression norms
through depth eight increase toward the predicted value.

For $d>1$, every output word in \eqref{eq:right} has at least $d-1$ more
ones than the input.  In particular, $T_f$ is not onto and cannot produce a
Haar-equivalent affine basis.  More importantly, terminal families from
different chaos orders are nested rather than suffix-free.  Their ranges
need not be orthogonal, and the Pythagorean reduction \eqref{eq:factor}
breaks down.  This is why the theorem is a proper slice of, not a full
solution to, the source's unrestricted commutant problem.

\section{Novelty and review status}

The earlier paper \textup{[2]} treats Haar chaoses of arbitrary fixed order
and characterizes simultaneous boundedness in every $L^p$, while its exact
continuous-invertibility result is for order one.  The present source proves
the single-$L^2$ scalar formula for order one.  A bounded search using the
exact titles together with fixed-chaos, commutant, and mixed-Hardy terms
found no explicit formula \eqref{eq:norm}.  Since the abstract
free-semigroup commutant model is standard, novelty is claimed only for the
explicit fixed-Haar-chaos reduction and exact mixed norm, pending specialist
review.

\medskip
\noindent\textbf{Human-review focus.}
Check the orientation in \eqref{eq:right}, the suffix-free argument, the
infinite-coordinate passage in Lemma~\ref{lem:column}, and the factor
$2^{-(d-1)/2}$ in \eqref{eq:norm}.

\medskip
\noindent{\large\bfseries References}\par\smallskip
\noindent[1] S.~V. Astashkin and P.~A. Terekhin,
\emph{Sequences of dilations and translations equivalent to the Haar system
in $L^p$-spaces}, arXiv:2005.04648.\par
\smallskip
\noindent[2] S.~V. Astashkin and P.~A. Terekhin,
\emph{Sequences of dilations and translations in function spaces},
J. Math. Anal. Appl. \textbf{457} (2018), 645--671.\par

\end{document}
